\chapter{Discussion et limites}
\label{chapitre:discussion}

\section{Réponse à la question principale}

Le résultat le plus défendable concerne HYPO. Une détection directionnelle, individualisée,
fondée sur plusieurs signaux et persistante produit des épisodes temporellement plus cohérents avec les
dégradations graduelles que les comparateurs IF et LOF. Après soustraction de l'exécution
propre, HYPO couvre 29 des 33 événements graduels et atteint IoU20 dans 13 cas. Cette preuve
reste technique: elle montre que le code réagit à des baisses contrôlées, non qu'il détecte
une boiterie clinique.

Le taux de nouveaux départs de 57,6\,\% nuance la couverture de 87,9\,\%. Une partie de la
réponse prolonge un épisode déjà actif au lieu d'en créer un. Le délai médian de 20,75 heures
montre également que le terme \textit{précoce} doit être interprété relativement à des
dégradations de 36 à 60 heures, et non comme une réaction immédiate.

\section{Niveau de preuve des principales conclusions}

Le Tableau~\ref{tab:portee_conclusions} associe à chaque conclusion son niveau de preuve.

\begin{table}[H]
  \centering
  \caption{Portée des conclusions principales.}
  \label{tab:portee_conclusions}
  \small
  \begin{tabular}{p{0.22\textwidth}p{0.34\textwidth}p{0.35\textwidth}}
    \toprule
    Niveau & Conclusion permise & Conclusion non permise \\
    \midrule
    Démontré techniquement
      & HYPO réagit à des baisses graduelles après la période de référence et localise mieux ces scénarios
        que les comparateurs IF/LOF.
      & HYPO constitue un détecteur clinique ou distingue la cause de la dégradation. \\
    Soutenu mais limité
      & La persistance, la direction et la cohérence entre signaux améliorent le compromis
        entre recouvrement et charge de fond dans cette campagne.
      & Les paramètres constituent un optimum physiologique ou sont transférables à une
        autre ferme sans réévaluation. \\
    Exploratoire
      & INSTABILITÉ peut produire une sortie de surveillance distincte et la fusion peut
        représenter une séquence synthétique.
      & L'instabilité précède habituellement une boiterie réelle. \\
    Non démontré
      & Aucune
      & Sensibilité, spécificité, valeur prédictive ou bénéfice clinique. \\
    \bottomrule
  \end{tabular}
\end{table}

\section{Apport de l'ablation et de l'OFAT}

L'ablation indique que la quantité de points anormaux n'est pas le critère décisif. IF
ponctuel produit 7,658 départs par vache-jour, mais aucun événement à IoU20. IF et LOF suivis
de persistance réduisent la charge, sans localiser les fenêtres avec IoU20. HYPO atteint un
IoU attribuable moyen de 0,137 et 29,5\,\% à IoU20 sur les quatre profils. La direction du
changement, la cohérence entre familles et l'accumulation temporelle constituent donc les
éléments les mieux soutenus par l'expérience.

L'analyse OFAT montre que les conclusions ne reposent pas sur une valeur unique de chaque
seuil. Elle identifie toutefois deux choix structurants: l'agrégation et le nombre de familles
concordantes. Exiger quatre familles réduit fortement la couverture et le fond; n'en exiger
que deux augmente la charge. Une agrégation courte réagit plus vite mais localise moins bien,
tandis qu'une agrégation longue perd des événements. Les valeurs de référence constituent un
compromis documenté, pas un optimum. Modifier les seuils après lecture de cette analyse aurait
transformé une étude de robustesse en sélection sur les données d'évaluation; aucun réglage
a posteriori n'a donc été effectué.

La tentative de calibration par la dispersion du score propre à chaque période de référence confirme
ce principe de prudence. Elle ne réduit ni le fond moyen ni sa variabilité, et diminue IoU20
de 39,4 à 33,3\,\%. Elle est donc rejetée. La courbe détection--charge montre par ailleurs
qu'un seuil de 0,14 paraît favorable sur cette campagne, mais l'adopter après observation
des mêmes événements introduirait un biais de sélection. Le résultat sert à formuler une
hypothèse de développement ultérieur, non à améliorer rétroactivement la performance
annoncée.

\section{Portée de l'extension bidirectionnelle}

L'extension répond à une question secondaire: peut-on représenter une instabilité courte sans
l'assimiler automatiquement à une alerte de vérification? L'architecture y répond sur le plan
logiciel grâce à deux persistances indépendantes et à une sortie de surveillance. Les
contrôles brefs et l'activité coordonnée de type œstrus sont rejetés dans la campagne
synthétique, et la séquence instabilité--hypoactivité est récupérée dans 90,9\,\% des cas.

Ces résultats ne démontrent toutefois pas que cette séquence existe fréquemment avant une
boiterie réelle. Les profils d'instabilité et leur chronologie sont des hypothèses de test.
La fusion hiérarchique augmente la couverture des scénarios de vérification de 61,4 à 72,7\,\% par rapport à
HYPO seul dans la campagne étendue, mais elle n'améliore pas IoU20 et augmente le fond de
0,446 à 0,571 notification par vache-jour. Elle ne domine donc pas HYPO. Son intérêt actuel
est de structurer une hypothèse et une interface de surveillance, non de remplacer la méthode
principale.

\section{Contrôles et facteurs confondants}

Le rejet des 33 contrôles brefs vérifie qu'un pic capteur, deux heures d'exercice synthétique
ou une manipulation d'une heure ne suffisent pas dans les conditions testées. Cette réussite
dépend de la représentation choisie et ne garantit pas le rejet de tous les événements réels
de même nature.

La confusion avec l'hypoactivité non locomotrice est la limite centrale. Neuf scénarios sur
onze sont alertés, car leur forme est volontairement identique à une dégradation locomotrice
sur les signaux disponibles. Aucun réglage temporel ne peut entièrement séparer deux causes
qui produisent les mêmes observations. La température, l'ingestion, la production laitière,
les traitements, le vêlage ou un examen clinique sont nécessaires pour améliorer la
spécificité causale.

\section{Charge opérationnelle et performance}

Le fond de 0,402 notification par vache-jour pour HYPO et de 0,571 pour la fusion
hiérarchique demeure élevé pour un grand troupeau. Ces valeurs ne sont pas des taux de faux
positifs: une donnée propre peut contenir une dégradation réelle non annotée. Elles indiquent
néanmoins une fatigue d'alerte potentielle. Un déploiement devrait regrouper les alertes par
épisode, limiter les répétitions, prioriser les trajectoires persistantes et enregistrer le
résultat de la vérification humaine.

Le temps médian de 10,64 secondes pour 28 vaches montre que le prototype est techniquement
exécutable sur le corpus étudié. IF représente la majeure partie du coût, tandis que le calcul
HYPO+INSTABILITÉ+fusion représente 6,4\,\% du temps instrumenté. Ce résultat ne préjuge pas de
la latence en ferme, du stockage continu ou de la concurrence entre utilisateurs.

\section{Portée de l'analyse SLS}

L'AUC exploratoire de 0,924 et le test de Mann--Whitney sont favorables au nombre de
notifications, mais seulement trois animaux ont un SLS~$\geq2$ et ils appartiennent tous au
groupe \textit{Exercise}. Le score maximal et la surveillance d'instabilité ne séparent pas
les groupes. L'association peut refléter l'exercice, une autre différence de groupe ou une
instabilité d'échantillonnage. L'analyse vérifie que la chaîne de traitement peut être appliquée avant un
score réel; elle ne permet ni de choisir les seuils ni d'estimer une sensibilité clinique.

\section{Mise en perspective avec la littérature}

Les résultats s'accordent avec la signature comportementale la mieux établie de la boiterie~:
une réduction de l'activité locomotrice et une augmentation du temps couché
\citep{ito2010lying,blackie2011lying,chapinal2010measurement}. La branche HYPO cible
précisément cette direction. Les revues soulignent toutefois que les variables utiles et
leurs amplitudes dépendent du capteur, du protocole et du contexte
\citep{oleary2020accelerometers,alsaaod2019automatic}. La référence propre à chaque vache est
donc un choix méthodologique prudent pour absorber une partie de l'hétérogénéité observée;
elle ne constitue pas, à elle seule, une preuve de spécificité clinique.

L'agitation, les reports de poids et les changements posturaux peuvent être observés lors
d'un inconfort, mais les agrégats IceQube ne mesurent ni l'appui relatif des membres ni la
démarche. Une hausse de mouvement ou de transitions reste en outre compatible avec l'œstrus,
l'exercice, une manipulation ou un artefact. C'est pourquoi la branche INSTABILITÉ filtre
l'activité coordonnée et demeure exploratoire. Enfin, les performances rapportées varient
fortement d'une étude à l'autre
\citep{oleary2020accelerometers,qiao2021review}, ce qui rappelle qu'aucune signature unique
n'est pleinement fiable sur ces capteurs et conforte le cadrage en alerte non diagnostique.

\section{Comparaison avec les systèmes commerciaux}

Plusieurs systèmes commerciaux fondés sur des podomètres et des accéléromètres d'activité
sont déployés en élevage. Leur logique interne et leurs protocoles d'évaluation ne sont pas
toujours décrits avec le niveau de détail nécessaire à une comparaison indépendante
\citep{rutten2013invited,stangaferro2016use}. Le présent travail se distingue par un code
auditable, des paramètres figés, un protocole explicite et un manifeste d'intégrité. Il ne
revendique aucune performance supérieure à ces systèmes, faute d'étalon et de protocole
communs.

\section{Implications pratiques pour l'éleveur}

À l'échelle d'un troupeau, le système se positionne comme un outil de priorisation, non
d'alarme. La charge de fond mesurée (0,402 notification par vache-jour pour HYPO) reste trop
élevée pour une vérification systématique~; un déploiement réaliste afficherait plutôt chaque
jour les quelques animaux les plus anormaux, à partir du classement de l'interface. Le
bénéfice attendu est d'orienter plus tôt l'attention humaine vers des vaches dont le
comportement dérive, le coût étant le temps de vérification et le risque de fatigue d'alerte.
Tant qu'aucune validation clinique n'établit une sensibilité réelle, ce bénéfice demeure une
aide à l'observation et non une substitution au jugement de l'éleveur ou du vétérinaire.

\section{Forces et limites}

Le Tableau~\ref{tab:forces_limites} synthétise les forces et les limites du système.

\begin{table}[H]
  \centering
  \caption{Forces et limites du travail.}
  \label{tab:forces_limites}
  \small
  \begin{tabular}{p{0.43\textwidth}p{0.48\textwidth}}
    \toprule
    Forces & Limites \\
    \midrule
    Séparation temporelle stricte et exécution propre de référence.
      & Absence de labels cliniques dans le corpus principal. \\
    Profils graduels, multivariés et physiquement contraints.
      & Amplitudes partiellement hypothétiques. \\
    Ablation sur les mêmes événements et agrégation par vache.
      & Onze vaches seulement dans les campagnes. \\
    OFAT de vingt variantes pour HYPO.
      & Aucune interaction entre paramètres ni jeu de développement indépendant. \\
    Contrôles brefs et facteurs confondants explicites.
      & Hypoactivité locomotrice et non locomotrice non identifiables avec ces seuls signaux. \\
    Évaluation du temps d'exécution de la chaîne complète et artefacts reproductibles.
      & Temps mesuré sur une seule machine et un seul corpus. \\
    Analyse SLS fondée sur des mesures antérieures au score.
      & Trois SLS~$\geq2$, tous dans le groupe \textit{Exercise}. \\
    \bottomrule
  \end{tabular}
\end{table}

\section{Menaces à la validité}

L'analyse suit la typologie classique des menaces à la validité en recherche empirique
\citep{shadish2002experimental,wohlin2012experimentation}.

\subsection{Validité de construit}
Les injections représentent des changements de comportement, non des cas de boiterie. Les
capteurs ne mesurent pas directement la douleur, la lésion ou l'asymétrie de démarche. Le
seuil IoU20 est une convention technique et ne possède pas de signification clinique.

\subsection{Validité interne}
La séparation chronologique et la soustraction de l'exécution propre réduisent la fuite et
l'attribution erronée. Plusieurs scénarios proviennent toutefois de la même vache; l'unité
indépendante reste donc la vache. La graine unique améliore la traçabilité, mais ne mesure pas
toute la variabilité possible du placement des événements.

\subsection{Validité externe}
Le corpus principal provient d'une ferme et d'une période. Les différences de logement, de
sol, de saison, de conduite et de capteur peuvent modifier les distributions. La cohorte SLS
provient du même site et ne constitue pas une validation inter-fermes.

\subsection{Validité statistique}
Les proportions synthétiques sont descriptives. Les tests d'ablation utilisent onze paires
au niveau vache, ce qui respecte la structure mais limite la puissance. Les analyses OFAT et
INSTABILITÉ examinent plusieurs variantes sans jeu de test indépendant. Les AUC et valeurs de
$p$ de l'analyse SLS sont fragiles en raison du petit effectif et de la multiplicité. La
calibration individuelle et la courbe détection--charge sont des analyses exploratoires
postérieures à la validation principale; aucune valeur observée dans ces analyses n'est
utilisée pour remplacer la configuration de référence.

\section{Usage prévu et implications éthiques}

L'usage défendable est un tri des animaux à observer. Une notification signifie que plusieurs
dimensions du comportement ont changé de manière persistante par rapport à sa période de
référence; elle ne signifie pas que la vache soit boiteuse. Une vérification doit inclure l'observation
locomotrice, l'examen du pied si nécessaire et une consultation vétérinaire selon le contexte.
L'absence d'alerte ne garantit pas l'absence de douleur et le système ne doit automatiser ni
traitement ni réforme.

La gouvernance des données mérite d'être explicitée en vue d'un déploiement. Les droits de
propriété et d'usage ne peuvent pas être présumés~: ils doivent être définis entre
l'exploitation, le fournisseur technologique et les chercheurs, avec des règles explicites
de consentement et de confidentialité \citep{wolfert2017big,carbonell2016ethics}. Dans ce
travail, l'extraction CowAlert demeure locale, est exclue des artefacts diffusables et ne sert
qu'au contrôle de provenance~; seules des sorties agrégées et leurs empreintes d'intégrité
sont conservées dans la chaîne reproductible. Un système déployé devrait documenter la
finalité, les droits d'accès, la sécurité, la durée de conservation et l'audit des données,
et garantir qu'une alerte n'entraîne aucune décision automatique de traitement ou de réforme.

\section{Travaux nécessaires avant une revendication clinique}

Une validation clinique devrait réunir des scores locomoteurs datés, plusieurs évaluateurs,
des examens des pieds ou diagnostics vétérinaires et les données capteurs des jours
précédents. Le développement et le test devraient être séparés par vache, période et,
idéalement, ferme. Des contrôles réels doivent inclure œstrus, exercice, manipulation,
vêlage, chaleur et maladies non locomotrices. Les seuils pourraient alors être optimisés sur
le jeu de développement, puis gelés avant le test.

\section{Bilan}

Le module HYPO constitue le résultat technique principal: il est mieux étayé que les
comparateurs ponctuels, mais sa localisation et sa charge restent imparfaites. L'extension
bidirectionnelle organise une hypothèse complémentaire et révèle ses propres compromis sans
masquer le résultat négatif sur les confondants. Cette séparation entre preuve principale et
exploration rend l'ensemble plus cohérent et évite de transformer une hypothèse comportementale
en affirmation clinique.
